% ============================================================
% Section 154: δ势与无限高势壁附近的束缚态
% ============================================================
\section{量子力学：$\delta$势与无限高势壁附近的束缚态}
\label{sec:delta-wall-bound-state}

\begin{modelbox}[题目：$\delta$势与无限高势壁附近的束缚态能量修正]
粒子在一维势场中运动，
\[
V(x)=\begin{cases}\infty,&x\leq-a\\-V_0\delta(x),&x>-a\end{cases}
\]
求出当势壁离粒子很远时，对束缚态能量的修正值，并说明多远才算是``很远''；至少存在一个束缚态时，$V_0$ 和 $a$ 应满足什么严格条件？
\end{modelbox}

\begin{tipbox}[考点快速分析]
\begin{itemize}
    \item \textbf{分区求解}：$x\leq-a$ 处 $\psi=0$；$x>-a$ 区含 $\delta$ 势
    \item \textbf{衰减常数}：$\kappa=\sqrt{-2mE}/\hbar$，特征长度 $b=\hbar^2/(mV_0)$
    \item \textbf{边界条件}：$x=-a$ 处 $\psi=0$；$x=0$ 处 $\psi$ 连续、$\psi'$ 跃变 $-2/b$
    \item \textbf{远离判据}：$\kappa a\gg1$，即 $a\gg b/2=\hbar^2/(2mV_0)$
\end{itemize}
\end{tipbox}

\begin{derivationbox}[标准解题过程]
\textbf{(1) 势壁远离时的能量修正}

设束缚态能量 $E<0$，衰减常数 $\kappa=\sqrt{-2mE}/\hbar$，特征长度 $b=\hbar^2/(mV_0)$。

\textit{分区波函数}：
\begin{itemize}
    \item 区域 I（$-a<x<0$）：$\psi_{\text{I}}(x)=Ae^{\kappa x}+Be^{-\kappa x}$
    \item 区域 II（$x>0$）：$\psi_{\text{II}}(x)=e^{-\kappa x}$（取 $C=1$）
\end{itemize}

\textit{边界条件}：
\begin{itemize}
    \item $x=-a$：$\psi(-a)=0\implies A=-Be^{2\kappa a}$
    \item $x=0$：$\psi$ 连续 $\implies A+B=1$；$\psi'$ 跃变 $\implies-\kappa-\kappa(A-B)=-2/b$
\end{itemize}

消去 $A$、$B$ 得能量本征方程：
\[
\kappa b=1-e^{-2\kappa a}.
\]

\textit{势壁远离近似}（$\kappa a\gg1$）：$e^{-2\kappa a}\ll1$，方程近似为
\[
\kappa\approx\frac{1}{b}\left(1-e^{-2\kappa_0 a}\right),\quad\kappa_0=\frac{1}{b}=\frac{mV_0}{\hbar^2}.
\]

零级近似（孤立 $\delta$ 势阱）：$E_0=-\dfrac{mV_0^2}{2\hbar^2}$。能量修正：
\begin{equation}
\boxed{\Delta E\approx\frac{mV_0^2}{\hbar^2}\exp\left(-\frac{2mV_0 a}{\hbar^2}\right).}
\end{equation}

``很远''的判据：$a\gg\dfrac{\hbar^2}{2mV_0}$，即势壁距离远大于孤立 $\delta$ 势阱中波函数的衰减长度。

\textbf{(2) 存在束缚态的严格条件}

令 $\varepsilon=\kappa b$，方程变为 $\varepsilon=1-e^{-(2a/b)\varepsilon}$。定义 $f(\varepsilon)=\varepsilon-1+e^{-(2a/b)\varepsilon}$。

$f(0)=0$，$f'(0)=1-\dfrac{2a}{b}$。

若 $f'(0)\geq0$（即 $b\geq2a$ 或 $V_0\leq\dfrac{\hbar^2}{2ma}$），函数在 $(0,\infty)$ 上恒正，无正根，\textbf{无束缚态}。

若 $f'(0)<0$（即 $b<2a$ 或 $V_0>\dfrac{\hbar^2}{2ma}$），函数从 0 下降后趋于 $+\infty$，必存在唯一正根，\textbf{有且仅有一个束缚态}。

\begin{equation}
\boxed{\text{存在束缚态的严格条件：}V_0>\frac{\hbar^2}{2ma}.}
\end{equation}
\end{derivationbox}

\begin{formulabox}[答案汇总]
\begin{center}
\begin{tabular}{ll}
\toprule
\textbf{结论} & \textbf{结果} \\
\midrule
能量修正 & $\Delta E\approx\dfrac{mV_0^2}{\hbar^2}\exp\left(-\dfrac{2mV_0 a}{\hbar^2}\right)$ \\
``很远''判据 & $a\gg\dfrac{\hbar^2}{2mV_0}$ \\
束缚态条件 & $V_0>\dfrac{\hbar^2}{2ma}$ \\
\bottomrule
\end{tabular}
\end{center}
\end{formulabox}

\begin{insightbox}[物理图像]
\begin{itemize}
    \item 无限高势壁对 $\delta$ 势阱的束缚态起``排斥''作用：将势壁移近（$a$ 减小）会升高能量，甚至使束缚态消失。
    \item 存在束缚态的条件 $V_0>\hbar^2/(2ma)$ 表明：势阱必须足够深，或势壁足够远，才能将粒子束缚在 $x=0$ 附近。
    \item 该模型是理解表面态、量子点受限效应等问题的简化原型。
\end{itemize}
\end{insightbox}
